% !TeX encoding = UTF-8
% !TeX root = thesis.tex

\section{Einleitung}

% ---------------------------------------------------------------------------
\subsection{Motivation}
% ---------------------------------------------------------------------------

Automatisierte Anomalie-Erkennung ist ein zentrales Werkzeug im Betrieb großer
Cloud-Infrastrukturen. Monitoring-Systeme werten kontinuierlich Systemmetriken aus
und lösen bei auffälligen Abweichungen Alerts aus. In der Praxis führt dieser Ansatz
jedoch zu einem hohen Aufkommen an Fehlalarmen (\emph{False Positives}):
Betriebsschwankungen, reguläre Lastspitzen oder geplante Batch-Jobs erzeugen
Metrikanomaliemuster, die denen echter Angriffe oder Systemfehler ähneln, ohne es zu sein.

Dieses Phänomen ist in der Literatur als \emph{Alert Fatigue} bekannt und stellt ein
anerkanntes operatives Problem im Sicherheitsbetrieb dar. Es beeinträchtigt die
Reaktionsfähigkeit von Analysten auf tatsächliche Vorfälle. Die Reduktion der
False-Positive-Rate ist daher eine zentrale Anforderung an praxistaugliche
Erkennungssysteme. Dabei soll die Erkennungsrate sicherheitsrelevanter Ereignisse
nicht wesentlich sinken. Abschnitt~2.1 ordnet diesen Befund in die Literatur ein.

% ---------------------------------------------------------------------------
\subsection{Problemstellung}
% ---------------------------------------------------------------------------

Ein grundlegendes Problem metrischer Anomalie-Detektoren besteht im Zielkonflikt
zwischen Recall und False-Positive-Rate. Um sicherzustellen, dass keine echten
Angriffe unentdeckt bleiben (hoher Recall), muss der Entscheidungs-Threshold des
Detektors niedrig gesetzt werden. Ein niedriger Threshold markiert jedoch
zwangsläufig auch viele tatsächlich unbedenkliche Instanzen als auffällig, was
die FPR in die Höhe treibt. Dieser Trade-off ist bei der Verwendung von wenigen,
universell verfügbaren Metriken besonders ausgeprägt, da solche Features allein
häufig keine ausreichende Trennschärfe zwischen malignen und benignen Fällen
bieten.

Log-Dateien liefern einen ergänzenden, qualitativ anderen Informationstyp: Sie
enthalten prozessuale Beschreibungen von Systemereignissen, die den Kontext einer
Metrik-Auffälligkeit klären können~\cite{he2021survey}. Eine Lastspitze in der
CPU-Auslastung kann sowohl bei einem malignen Prozess als auch bei einem legitimen
Backup-Job auftreten. Der zugehörige Log-Eintrag hingegen unterscheidet beide
Fälle in der Regel klar. Log-Daten eignen sich daher als nachgelagerte
Verifikationsstufe: Alerts, die metrisch ausgelöst wurden, können auf Case-Ebene
anhand ihres Log-Kontexts gefiltert werden.

Die Kombination beider Signalquellen zu einer hybriden Pipeline ist konzeptuell
naheliegend. Die Pipeline besteht aus einem metrischen Trigger, gefolgt von einer
log-basierten Verifikation. CloudAnoBench~\cite{zou2025cloudanobench} stellt einen
synthetischen Benchmark bereit, der für jeden überwachten Case sowohl Metriken als
auch Log-Dateien bereitstellt und damit eine kontrollierte Evaluation eines solchen
zweistufigen Ansatzes ermöglicht.

% ---------------------------------------------------------------------------
\subsection{Zielsetzung und Forschungsfrage}
% ---------------------------------------------------------------------------

Ziel dieser Arbeit ist die Entwicklung und Evaluation einer zweistufigen hybriden
Anomalie-Erkennungspipeline auf CloudAnoBench. Der Schwerpunkt liegt auf einer
methodisch kontrollierten Evaluation und der expliziten Diskussion der Grenzen des
Ansatzes, einschließlich unerwarteter Befunde.

Die zentrale Forschungsfrage lautet:

\begin{quote}
  \emph{Kann eine zweistufige hybride Pipeline, bestehend aus einem metrischen
  Anomalie-Detektor und einer log-basierten Nachfilterung, die
  False-Positive-Rate auf dem CloudAnoBench-Datensatz deutlich senken, ohne die
  Recall-Rate unter 0{,}95 zu drücken?}
\end{quote}

Die Arbeit untersucht diese Frage durch einen kontrollierten Ablations-Aufbau, der
eine rein metrische Baseline, eine log-basierte Komponente sowie zwei
Fusionsstrategien miteinander vergleicht, ein binäres AND-Gate und eine Soft-Fusion. Die Evaluation folgt dabei strikten Anforderungen an die Vermeidung
von Datenleck und eine leakage-sichere Trainings-Test-Trennung.

Die Skalierbarkeitsbetrachtung dieser Arbeit ist als konzeptionelle Einordnung
der modularen und grundsätzlich parallelisierbaren Pipeline zu verstehen.
Eine empirische Untersuchung der Skalierbarkeit durch Lasttests, Durchsatz- oder
Latenzmessungen beziehungsweise eine verteilte Implementierung ist nicht
Gegenstand der Arbeit (vgl.\ Abschnitt~\ref{sec:scalability}).

% ---------------------------------------------------------------------------
\subsection{Beiträge der Arbeit}
% ---------------------------------------------------------------------------

Die vorliegende Arbeit leistet die folgenden konkreten Beiträge:

\begin{enumerate}
  \item \textbf{Reproduzierbare Case-Level-Pipeline auf CloudAnoBench.} Es wird eine
    in Notebooks dokumentierte, reproduzierbare Pipeline entwickelt. Diese deckt alle
    Phasen ab, von der Datenaufbereitung über Case-Level-Feature-Engineering und
    Modelltraining bis hin zu Fusion und Evaluation auf CloudAnoBench. 67~automatisierte Tests in
    14~Testklassen sichern zentrale Implementierungs- und Konsistenzbedingungen ab;
    die Reproduzierbarkeit wird zusätzlich durch fixierte Seeds, gespeicherte
    Split-Zuordnungen, die dokumentierte Ausführungsreihenfolge und die
    spezifizierte Umgebung unterstützt. Der finale
    Testsplit wird mit \texttt{random\_state=42} fixiert; der interne Validierungsanteil
    wird ohne Zugriff auf das Testset über einen festen Seed aus einer vordefinierten
    Kandidatenliste gewählt.

  \item \textbf{Case-Level-Metrik-Baseline mit XGBoost auf aggregierten Features.}
    Es wird eine metrische Baseline mit XGBoost~\cite{chen2016xgboost} entwickelt,
    die direkt auf 25~aggregierten Case-Level-Features der fünf universellen Systemmetriken
    trainiert wird (Mittelwert, Maximum, Standardabweichung, 95.~Perzentil und Trend
    je Metrik). Die Beschränkung auf die fünf universell verfügbaren Features ist eine
    methodisch begründete Entscheidung zur Vermeidung von Target Leakage durch
    strukturelle Fehlwerte.

  \item \textbf{Log-basierte Verifikation mit Normalisierung, TF-IDF und logistischer
    Regression.} Es wird eine Case-Level-Komponente implementiert, die Log-Texte
    vor der TF-IDF-Vektorisierung durch eine regelbasierte Normalisierung synthetischer
    Artefakte bereinigt (Zeitstempel, IP-Adressen, PIDs, Ports, isolierte Zahlen werden
    durch generische Token ersetzt). Die logistische Regression auf dem normalisierten
    TF-IDF-Raum wird als interpretierbare Baseline eingesetzt; das Vokabular wird
    ausschließlich auf Trainingsdaten gefittet.

  \item \textbf{Ablation der Fusionsstrategien einschließlich Degenerationsbefund.}
    Es werden zwei Fusionsstrategien evaluiert, AND-Gate und Soft-Fusion,
    und mit den Einzelkomponenten verglichen. Ein zentraler empirischer Befund
    ist die vollständige Degeneration des AND-Gates: Der metrische Detektor löst
    beim gewählten Threshold alle 253~Test-Cases aus (100~\%), sodass das AND-Gate
    funktional äquivalent zur reinen Log-Filterung wird. Soft-Fusion verschlechtert
    sich gegenüber Log-only. Dieser Befund wird in Abschnitt~\ref{sec:degeneration}
    explizit diskutiert.

  \item \textbf{Szenarioaufgelöste Fehleranalyse.}
    Es wird eine diagnostische Case-Level-Feh\-ler\-analyse durchgeführt, die
    False Negatives und False Positives nach Szenario aufschlüsselt, ohne
    Modelle neu zu trainieren oder Thresholds zu verändern. Die Analyse
    lokalisiert die Fehlerquellen und stützt die Interpretation der Hauptergebnisse
    (vgl.\ Abschnitt~\ref{sec:error_analysis}).
\end{enumerate}

% ---------------------------------------------------------------------------
\subsection{Aufbau der Arbeit}
% ---------------------------------------------------------------------------

Abschnitt~2 legt die konzeptionellen und methodischen Grundlagen der Pipeline dar.
Behandelt werden Cloud-Anoma\-lie\-erkennung, unbalancierte Klassifikation und
Precision-Recall-Metriken, XGBoost, TF-IDF, Data Leakage und gruppenbasierte
Validierung sowie Klassifikatorfusion und Case-Level-Evaluation.

Abschnitt~\ref{sec:data_preparation} beschreibt CloudAnoBench sowie die Daten\-aufbereitung
und leakage-sichere Split-Strategie. Abschnitt~4 stellt die metrische Baseline
(XGBoost auf aggregierten Case-Level-Features) vor; Abschnitt~5 beschreibt die
log-basierte Verifikationskomponente (TF-IDF und logistische Regression).
Abschnitt~6 präsentiert die hybride Fusion und die Ablationsstudie einschließlich
der AND-Gate-Degeneration sowie der szenarioaufgelösten Fehleranalyse. Die
Ergebnisse werden durch Tabellen und Abbildungen dokumentiert.

Abschnitt~7 diskutiert die Ergebnisse im Hinblick auf die Forschungsfrage,
bewertet die Rolle der einzelnen Komponenten und benennt die methodischen
Limitationen der Arbeit. Abschnitt~8 fasst die Befunde zusammen und gibt einen
Ausblick auf mögliche Erweiterungen.
